\subsection{Finite checkpoint selection and the shortest native two-record order}
\label{sec:source-checkpoint-selection}

The preceding continuation theorem can be strengthened from an occurrence
bound under supplied IID proposals to an explicit finite information
projection and its causal implementation. The interface remains declared:
the public record map, preparation domain, allowed native moves, retained
samples, finite deadline, full-history entropy cover and reference are
inputs. A1's availability of a checkpoint does not imply mandatory
publication of this record by this deadline. This section changes the
feasible family by imposing that requirement; it does not call the result
the original unconstrained proposal law or a full A1--A3 derivation of M1.

\paragraph{Complete record constraints and finite selection.}
For a public map $r:X\to R$ and observation $o:X\to O$, agreement of
every test $g(r)$ is equivalent to
$o(x)=o(y)\Rightarrow r(x)=r(y)$. Equality predicates prove the
reverse implication. For a tuple of records, testing all primitive
coordinates is therefore complete for its generated test algebra, including
nonlinear tests. This is not a completeness assertion about the entire A1
operational algebra. Preservation by the joint old-record/current-state
map is necessary and sufficient for bounded native continuation on connected
support, provided the current receiver scalar is determined by the retained
record. The topology-derived completion supplies sufficiency even after
irrelevant preparation distinctions have been lost.

At horizon $T$, let $\mathcal W_T$ contain every word in the declared
finite move alphabet and let $G_T$ be the words whose native observations
determine $r$. On the full simplex supported on $G_T$, a faithful inherited
reference $R_T$ has the unique information projection
\[
 Q_T(w)=\frac{R_T(w)\mathbf 1_{G_T}(w)}{Z_T},
 \qquad Z_T=\sum_{w\in G_T}R_T(w)>0.
\]
For every feasible $p$, its entropy gap is
\[
 D(p\Vert R_T)=-\log Z_T+
 \sum_{w\in G_T}Q_T(w)
 \left[\frac{p(w)}{Q_T(w)}\log\frac{p(w)}{Q_T(w)}
       +1-\frac{p(w)}{Q_T(w)}\right].
\]
Nonnegativity and the strict equality case prove uniqueness, including
candidate laws with zeros. The information cost is $-\log Z_T$.

For positive reference move weights $a(c,e)$, backward continuation masses
and the resulting online transition law are
\[
 H_0(c)=\mathbf1_{\{c\text{ publishes }r\}},\qquad
 H_{t+1}(c)=\sum_e a(c,e)H_t(A_ec),\qquad
 K_t(c,e)=\frac{a(c,e)H_{t-1}(A_ec)}{H_t(c)}.
\]
Here $c$ is a response checkpoint over all admissible preparations, so
neither an unknown payload nor a future sample is read by the controller.
The recursion sums every native continuation exactly once. Its positive
mass is equivalent to existence of a publishing continuation; connected
topology supplies a finite feasible deadline. Products of the normalized
transition probabilities telescope to $Q_T$. Every positive selected
history publishes, and every good history retains positive probability.
There is no retrospective rejection of failed physical runs. Exact
planning, retained records, random choice and decoding remain host
interfaces whose physical implementation cost is unsupplied.

\paragraph{A sharp path theorem.}
For the isolated path $0,1,\ldots,d$, $d\geq1$, put independent unknown
deviations $x,y$ at ports $0,1$, prepare all other deviations at zero, and
sample receiver $d$ initially and after every mean. Requiring both records
forces at least $2d-1$ means. Edge $0$ must transmit $x$. Each other cut
separates both unknowns from the receiver and needs at least two crossings:
one crossing supplies at most one independent scalar to the initially
calibrated downstream region. A successful word at equality therefore uses
edge $0$ once and each other edge twice.

Label the first crossings $F_1,\ldots,F_{d-1}$ and the second wave
$S_0,\ldots,S_{d-1}$. The successful minimal words are exactly the linear
extensions of
\[
 F_1<\cdots<F_{d-1},\qquad S_0<\cdots<S_{d-1},
 \qquad F_{j+1}<S_j\quad(0\leq j\leq d-2).
\]
A first crossing out of order averages zero deviations and wastes one of
the two required transmissions. If $S_j$ precedes $F_{j+1}$, the cut behind
it is exhausted while the downstream region holds only one source scalar;
for $j=0$, this is immediate erasure of the independent record difference.
The second crossings must occur in outward order to convey $x$ before each
cut is exhausted. These arguments give necessity. For sufficiency,
incomparable events have disjoint edges and their means commute, so every
linear extension has the final state of the serial two-wave word. Its two
informative receiver samples are
\[
 u=\frac{y}{2^{d-1}},\qquad
 v=\frac{x}{2^d}+\frac{(d+1)y}{2^{d+1}}.
\]
They identify both records. At $d=1$, $u$ is the initial receiver sample.
Deleting the final second-wave event gives Dyck words with $d-1$ events in
each wave, hence exactly $C_{d-1}=\binom{2d-2}{d-1}/d$ histories.

Every positive time-homogeneous edge-product reference assigns each of
these shortest words the same numerator
$a_0\prod_{i=1}^{d-1}a_i^2$. Thus its selected law is uniform over the
Catalan histories, independently of those edge weights. This does not
extend automatically to longer deadlines or arbitrary A3 reference rules.
The decoder is
\[
 y=2^{d-1}u,\qquad x=2^dv-(d+1)2^{d-2}u.
\]
Its sharp independent sample-error gains are $2^{d-1}$ and
$(d+5)2^{d-2}$. The linear operation count therefore coexists with
exponential precision amplification; no physical precision claim follows.

\paragraph{Evidence and scope.}
The finite grammar, continuation, entropy and online-policy theorems have
kernel-checked proofs, as do native commutation and the decoder and error
identities. The audit checks their transitive axioms. The full path-language
classification and Catalan necessity proof above are analytic, not
kernel-formalized. The independent executable packet checks every
shortest-horizon word through $d=5$ (1,969,761 words), and all 6,918 positive
histories through $d=10$, using native scalar basis replay. Producer Dyck
enumeration and verifier prerequisite-poset enumeration are independent.

On the captured seams $0,1,14,23,45$, both records require seven means and
exactly five shortest histories exist. All selected histories execute
seven means and retain eight samples. The sum alone takes four means and
does not determine either record. At eight means, 104 of 65,536 reference
words succeed and the first-move probabilities are $(0,94,5,5)/104$; the
seven-mean law starts on edge $1$ with certainty. Independently selected
deadline optimizers therefore need not form a compatible temporal family.
The control packet also retains tilted references, every signed
intervention and consumed writer in thirty online trajectories, complete
reference denominators and all selected native work. Every captured
preparation charges 15,360 writes. Expanded enumeration tapes are
recomputed from compact commitments.

Selecting the public record interface and metric menu, compatible
refinement and geometric identification, physical storage/control/precision,
and the immutable-version semantic event-order bridge remain separate
source obligations. This theorem does not select $L/\sqrt q$ or provide a
physical spacetime identification. Its finite temporal selection replaces a
supplied successful schedule inside the stated interface.
